\section{Introduction}\label{sec:intro}

Mountain glaciers are sensitive recorders of climate and a critical component of regional water resources, yet projecting their evolution requires computing their distributed surface mass balance (SMB)---the spatial pattern of accumulation and ablation---which remains poorly constrained at the global scale, particularly in remote parts of the world.
This is in part because the direct measurement of SMB relies on stake networks and snow pits that are labour-intensive and hazardous in steep accumulation areas. They are consequently deployed on less than 0.2\,\% of the world's glacierised area \citep{cogley2011}. The resulting records, compiled by the World Glacier Monitoring Service \citep{wgms2025}, cover only a few accessible benchmark glaciers \citep{zemp2019}.
Large-scale glacier models therefore rarely calibrate on stake data, which serve mainly as an independent check; instead, the parameters of their SMB models are tuned against glacier-wide geodetic mass balances \citep{marzeion2012, radic2014, hugonnet2021}.
As remote-sensing products of glacier surface
velocity \citep{millan2022} and
elevation change \citep{hugonnet2021, beraud2023} reach near-global coverage, the
question is no longer whether the data exist, but how to combine them to infer
physically consistent SMB fields.

\subsection{Existing approaches}\label{sec:sota}
SMB is most commonly obtained from glacier models that quantify the accumulation and ablation terms from meteorological inputs---at least precipitation and temperature, and further variables where a surface energy balance is resolved \citep{marzeion2012, rounce2020, buri2023}. These inputs come from climate models and reanalyses whose resolution is too coarse to represent high-mountain precipitation and temperature accurately, so the forcing must be bias-corrected. With calibration limited to glacier-wide geodetic balances, the problem stays strongly underdetermined: different parameter sets and model formulations reproduce the same integrated balance while implying different distributed SMB patterns and diverging projections \citep{rounce2020quantifying, schuster2023glacier}. 

Assimilation frameworks tackle this by calibrating against distributed observations rather than a glacier-wide total: FROST \citep{herrmann2025} tunes an elevation-dependent SMB model coupled to IGM ice dynamics to infer equilibrium-line altitudes across the Alps, fitting the parameters of a prescribed SMB form with an ensemble Kalman filter. Being derivative-free, it obtains parameter sensitivities from the spread of an ensemble that must grow with the number of unknowns, whereas differentiating end to end through our forward model returns the sensitivity to every parameter at roughly the cost of one additional forward run.
Alternatively, sparse direct measurements are extended to distributed fields by homogenising them against decadal geodetic volume change, as in the long-running Glacier Monitoring Switzerland programme \citep[GLAMOS;][]{glamos}. These provide a valuable observational reference, but rest on a large number of statistical assumptions and, more importantly, exist for only a handful of glaciers.

Satellite-derived velocity and elevation-change fields are now widely available. This has driven a family of approaches, closest to the present work, that infer the SMB from mass-conservation principles. The governing equation relates the rate of surface-elevation change to the divergence of the ice flux and the SMB \citep{cuffey2010}: given ice thickness and surface velocity, the vertically integrated flux can be computed, differentiated spatially, and combined with an observed elevation-change field to provide an SMB estimate. From early flux-gate estimates of emergence velocity in the ablation zone \citep{berthier2012, brun2018, kneib2022}, the approach has been generalised to fully distributed SMB over entire glaciers \citep{bisset2020, miles2021, pelto2021, vantricht2021}, to seasonal elevation-change time series \citep{zeller2023} and to debris-covered glaciers \citep{bhushan2024}. A recent realisation is that of \citet{kneib2024}, who reconstructed a high-resolution distributed SMB for Argentière Glacier (French Alps) over 2012--2021 from Pléiades elevation change and velocity, ground-penetrating-radar thickness and three different ice-thickness models.

Two difficulties beset this route. First, the flux divergence is the spatial derivative of the product of two independently measured noisy fields---surface velocity and ice thickness---and is therefore dominated by small-scale noise, which is suppressed by \emph{post hoc} spatial filtering \citep{vantricht2021, bisset2020}. Filtering also serves a physical purpose, since averaging over a few ice thicknesses represents the longitudinal stress coupling that governs how ice flow responds to variations in thickness and surface slope \citep{kamb1986, enderlin2016}; yet smoothing a quantity that respects mass conservation locally redistributes mass between grid cells and violates the very balance the method enforces. Per-glacier tuning of smoothing length scales, density assumptions and weighting schemes further hampers transfer to the regional scales most relevant for water-resource applications.
Second, the flux divergence is computed at a single observational snapshot and assumed representative of the whole observation period, whereas glacier geometry, velocity and flux all evolve under ongoing climate forcing \citep{johannesson1989, zekollari2015}; treating the flux as stationary introduces a bias that grows with the length of the window---a limitation \citet{kneib2024} acknowledge explicitly. Transient assimilation frameworks confront this directly, notably the dynamical spin-up employed by \citet{hartl2025} and the AGILE framework of \citet{schmitt2026}, which fit an evolving glacier model to observations rather than to a single snapshot.

\subsection{This study}\label{sec:contribution}
In this study we infer the SMB by inverting a transient glacier-evolution model: starting from a dynamically calibrated state, the mass-conservation equation is integrated forward from one DEM to the next, and reverse-mode automatic differentiation through the whole simulation returns the elevation-dependent SMB profile whose modelled geometry change best reproduces the observed one. Inversions of transient ice-flow models are not themselves new---\citet{brinkerhoff2025} jointly infer bed elevation, basal traction and surface mass balance for Sít' Tlein, with Bayesian quantification of the resulting uncertainty---but they draw on rich, glacier-specific observations; our aim here is to recover the SMB from two DEMs and a single velocity field alone.

Two features distinguish our formulation from flux-divergence inversions. First, the SMB is recovered by solving the forward mass-conservation equation rather than by differentiating observed fields: the flux divergence is recomputed at every time step from the evolving glacier state, so the inferred SMB is dynamically
consistent over the whole window, rather than assuming a stationary flux, and mass is conserved at every step, with no \emph{post hoc} smoothing of a noisy divergence field. Second, we parametrise the SMB as a one-dimensional function of surface elevation, evaluated at every grid cell to yield a 2D field;
elevation carries the first-order control on alpine SMB, so this leaves roughly ten unknowns in place of one per pixel and constrains an inversion that would otherwise be underdetermined, with no length scale or density assumption left to fix per site (Sect.~\ref{sec:inference}). The price is that horizontal variability within an elevation band, such as the avalanche-fed hotspots resolved by \citet{kneib2024}, is not recovered: a deliberate trade-off of sub-elevation detail for a transferable, tuning-free inversion, to which we return in Sect.~\ref{sec:discussion}.

The accuracy of our geometry-based inversion is ultimately controlled by that of the input DEMs. Near-global elevation-change products \citep{hugonnet2021} make the method broadly applicable, but their inaccuracies limit the fidelity of the recovered SMB, a penalty we quantify directly by repeating one of the inversions on such a product (Appendix~\ref{app:hugonnet}). Dedicated multi-temporal DEMs deliver the highest accuracy, and the rapidly expanding archive of such data is what makes a regional-scale application possible. We validate the method with two-epoch DEMs on three well-monitored Alpine glaciers spanning $\approx$12--80~km$^2$ in area---Argentière (French Alps), Great Aletsch and Rhône (Swiss Alps)---all inferred by the \emph{same} two-step pipeline, without per-glacier retuning: the basal-motion parameter is reduced to a single scalar fixed by an ensemble search on the surface-velocity misfit, and no thickness survey is assumed.

Sect.~\ref{sec:methods} details the two-step pipeline---the input and validation data (Sect.~\ref{sec:data}), data assimilation of the initial dynamical state, then SMB inference by
automatic differentiation through the transient forward run; Sect.~\ref{sec:results} presents the inferred SMB profiles and their validation against in situ stake measurements; and Sect.~\ref{sec:discussion} discusses the implications for regional-scale SMB inversion.
